\clearpage
\thispagestyle{empty}
\pdfbookmark[0]{General Theorems on Forces of Attraction and Repulsion. Edition and sources}{lehrsaetze-source-notice}
\begin{center}\Large\textbf{General Theorems on Forces of Attraction and Repulsion}\end{center}

AI-translated from the German text in Carl Friedrich Gauss, \textit{Werke}, volume V (Göttingen, 1867), pp.~195--242, with AI-produced typesetting. The edition contains the formal title and historical imprint, the repeated heading, all 37 articles and Gauss's note to Article 12. The original paper appeared in \textit{Resultate aus den Beobachtungen des magnetischen Vereins im Jahre 1839}, edited by Gauss and Weber (Leipzig, 1840).

An English version in Richard Taylor (ed.), \textit{Scientific Memoirs}, volume III, part X, Article VII (1843), pp.~153--196, was used for comparison. A.~K.~T.~Assis's edited and annotated opening excerpt, posted in October 2024, was also consulted. For modern German terminology, the comparison included Alexander Mirlin's \textit{Skript Elektrodynamik}, Chapter I, \textit{Elektrostatik}, at KIT.

\noindent\href{https://gdz.sub.uni-goettingen.de/id/PPN236006339}{German source (1867)}\quad
\href{https://archive.org/details/scientificmemoir03memo/page/153/mode/1up}{English version (1843)}\quad
\href{https://www.ifi.unicamp.br/~assis/Gauss-1840.pdf}{Assis's opening excerpt (2024)}\quad
\href{https://www.tkm.kit.edu/downloads/ws2012_theoc/Kapitel-I-Elektrostatik.pdf}{Mirlin, Elektrostatik}
\clearpage
